% Generated draft: requires mathematical and source review.
\documentclass[11pt]{article}
\usepackage[margin=25mm]{geometry}
\usepackage{fontspec}
\setmainfont{Latin Modern Roman}
\usepackage{amsmath,amssymb,mathtools,unicode-math}
\setmathfont{Latin Modern Math}
\usepackage{enumitem,needspace,xurl,hyperref}
\setlength{\parindent}{0pt}
\setlength{\parskip}{5pt}
\setlength{\emergencystretch}{3em}
\newcommand{\N}{\mathbb{N}}
\newcommand{\Z}{\mathbb{Z}}
\newcommand{\Q}{\mathbb{Q}}
\newcommand{\R}{\mathbb{R}}
\newcommand{\C}{\mathbb{C}}
\newcommand{\F}{\mathbb{F}}
\newcommand{\A}{\mathbb{A}}
\newcommand{\PP}{\mathbb{P}}
\newcommand{\E}{\mathbb{E}}
\newcommand{\eps}{\varepsilon}
\newcommand{\dd}{\,\mathrm{d}}
\begin{document}
\section*{UnsolvedMath ID: 701498}
\textbf{Source problem code:} catalogue entry
\subsection{Definitions and mathematical statement}
A variety of groups is a class of groups defined by group identities. A group \(G\) is relatively free if there exist a variety \(\mathcal V\) and a subset \(X\subseteq G\) generating \(G\) such that \(G\in\mathcal V\) and every map from \(X\) to any group \(H\in\mathcal V\) extends uniquely to a group homomorphism \(G\to H\). No restriction on the cardinality of \(X\) is specified in the source.

Interpret ``containing a free subsemigroup'' in its usual noncommutative sense: there exist \(a,b\in G\) such that evaluation at \(a,b\) is injective on the set of all nonempty words in two formal letters. Thus distinct nonempty positive words in \(a,b\) represent distinct elements of \(G\). The source does not explicitly specify the rank of the free subsemigroup; this interpretation requires confirmation.

Define the commutator subgroup by
\[
[G,G]=\langle g^{-1}h^{-1}gh\mid g,h\in G\rangle.
\]
The question is whether there exist a relatively free group \(G\), elements \(a,b\in G\) generating a free subsemigroup as above, an integer \(m\geq 0\), and elements \(c_1,\ldots,c_m\in[G,G]\) such that
\[
[G,G]=\langle c_1,\ldots,c_m\rangle.
\]
Here finite generation means generation as an abstract group.

A group is locally graded if every nontrivial finitely generated subgroup has a proper subgroup of finite index. The source additionally states that any group satisfying the requested conditions cannot be locally graded, citing the journal reference listed below. This is retained as a supplied necessary-condition assertion, not an independently verified result.
\subsection{Short English statement}
Does there exist a relatively free group containing a free semigroup on two generators whose commutator subgroup is finitely generated as a group? The source states that such a group would necessarily fail to be locally graded. The two-generator interpretation of ``free subsemigroup'' is explicit here but is not specified in the source.
\subsection{Sources}
\begin{itemize}
\item Local UnsolvedMath registry, permanent ID \(701498\), source problem code \texttt{catalogue entry}, titled ``Kourovka Notebook Problem 18.69.'' The supplied registry labels the problem open; this status has not been independently verified.
\item \emph{The Kourovka Notebook (Unsolved Problems in Group Theory)}, original problem 18.69, 18th issue (2014); the supplied provenance also identifies the 21st issue edition (2026), \texttt{arXiv:1401.0300v48}. Proposed by O. Macedo\'nska in 2014. The supplied collection description names E. I. Khukhro and V. D. Mazurov as editors.
\item The supplied statement cites \emph{Publ. Math. Debrecen} \textbf{81}, nos. 3--4 (2012), 415--420, for the assertion that the requested group cannot be locally graded. Article authors and title are not supplied.
\end{itemize}
\end{document}
